% courses/MetNum2023/Clases/assets/tikzstyles.tex
% ---------------------------------------------------------------------------
% Estilos compartidos para los diagramas del curso Métodos Numéricos
% (MetNum2023). Fuente única del "look" visual (colores, grosores, escalas)
% para que todas las figuras (matrices triangulares, curvas de convergencia
% / error, splines a trozos) sean consistentes entre clases.
%
% Uso: la clase debe cargar antes `tikz` y `pgfplots` (bloque §6.3 del
%      CLAUDE.md del curso), y luego
%      % courses/MetNum2023/Clases/assets/tikzstyles.tex
% ---------------------------------------------------------------------------
% Estilos compartidos para los diagramas del curso Métodos Numéricos
% (MetNum2023). Fuente única del "look" visual (colores, grosores, escalas)
% para que todas las figuras (matrices triangulares, curvas de convergencia
% / error, splines a trozos) sean consistentes entre clases.
%
% Uso: la clase debe cargar antes `tikz` y `pgfplots` (bloque §6.3 del
%      CLAUDE.md del curso), y luego
%      % courses/MetNum2023/Clases/assets/tikzstyles.tex
% ---------------------------------------------------------------------------
% Estilos compartidos para los diagramas del curso Métodos Numéricos
% (MetNum2023). Fuente única del "look" visual (colores, grosores, escalas)
% para que todas las figuras (matrices triangulares, curvas de convergencia
% / error, splines a trozos) sean consistentes entre clases.
%
% Uso: la clase debe cargar antes `tikz` y `pgfplots` (bloque §6.3 del
%      CLAUDE.md del curso), y luego
%      % courses/MetNum2023/Clases/assets/tikzstyles.tex
% ---------------------------------------------------------------------------
% Estilos compartidos para los diagramas del curso Métodos Numéricos
% (MetNum2023). Fuente única del "look" visual (colores, grosores, escalas)
% para que todas las figuras (matrices triangulares, curvas de convergencia
% / error, splines a trozos) sean consistentes entre clases.
%
% Uso: la clase debe cargar antes `tikz` y `pgfplots` (bloque §6.3 del
%      CLAUDE.md del curso), y luego
%      \input{../assets/tikzstyles.tex}
% (compilar desde el directorio de la clase para que la ruta relativa resuelva).
%
% Paleta propia de este curso (no se reusa la de Fisica3-2015/CDIV2017/
% ElecMag2024): mismos tonos base para alinear con keybox/notebox/warnbox,
% pero archivo independiente, sin dependencia cruzada entre cursos.
% ---------------------------------------------------------------------------

% ---- Paleta (alineada con las cajas del preámbulo: keybox/notebox/warnbox) ----
\definecolor{figblue}{HTML}{1F4E79}   % azul principal (L, curvas principales, keybox)
\definecolor{figred}{HTML}{B03A2E}    % rojo de énfasis (U, warnbox, error)
\definecolor{figamber}{HTML}{B8860B}  % ámbar (notebox, pivotes)
\definecolor{figgray}{HTML}{555555}   % auxiliares, ejes, ceros de matriz

% ---- TikZ: librerías y estilos reutilizables ----
% Nota: las puntas se declaran como `-{Latex[...]}` y no como `->`. La forma
% con llaves NO contiene un `>` literal, así que es inmune al `>` activo que
% introduce babel español (de todos modos se carga \usetikzlibrary{babel} en
% el bloque §6.3 para cubrir cualquier `->` suelto).
\usetikzlibrary{arrows.meta,calc,matrix,patterns}

\tikzset{
  figline/.style={figblue, line width=0.9pt},
  figaux/.style={figgray, dashed, line width=0.6pt},
  figaxis/.style={figgray, line width=0.7pt, -{Latex[length=2.2mm,width=1.5mm]}},
  % estructura de matrices (L triangular inferior, U triangular superior)
  matcelda/.style={minimum width=7mm, minimum height=7mm, anchor=center, font=\small},
  matnz/.style={matcelda, fill=figblue!12},
  matnzU/.style={matcelda, fill=figred!12},
  matcero/.style={matcelda, text=figgray!70},
  matpivote/.style={matcelda, fill=figamber!25, draw=figamber, line width=0.8pt},
  % nodos de interpolación / splines
  fignodo/.style={circle, inner sep=0pt, minimum size=4.5pt, draw=figblue, line width=0.9pt, fill=white},
  fignodoactivo/.style={fignodo, fill=figamber},
  % etiquetas
  figlbl/.style={font=\small},
  figlblaux/.style={font=\footnotesize, figgray},
}

% ---- pgfplots: estilos base ----
\pgfplotsset{
  % convergencia / error vs. n (grado, iteración) — suele ir en escala log
  errorfig/.style={
    width=10cm, height=6.5cm,
    xlabel style={font=\small},
    ylabel style={font=\small},
    tick label style={font=\footnotesize},
    every axis plot/.append style={line width=1.1pt, mark=*, mark size=1.5pt},
    grid=major, grid style={figgray!15},
    legend cell align=left,
    legend style={font=\footnotesize, draw=figgray!40, at={(0.5,-0.25)}, anchor=north},
    clip=false,
  },
  % interpolación / splines: un color por tramo, nodos marcados aparte
  interpfig/.style={
    width=10cm, height=6cm,
    axis lines=middle,
    xlabel style={font=\small, at={(axis description cs:1,0.5)}, anchor=west},
    ylabel style={font=\small, at={(axis description cs:0.5,1)}, anchor=south},
    tick label style={font=\footnotesize},
    every axis plot/.append style={line width=1.1pt},
    grid=none,
    legend cell align=left,
    legend style={font=\footnotesize, draw=figgray!40},
    clip=false,
  },
}

% (compilar desde el directorio de la clase para que la ruta relativa resuelva).
%
% Paleta propia de este curso (no se reusa la de Fisica3-2015/CDIV2017/
% ElecMag2024): mismos tonos base para alinear con keybox/notebox/warnbox,
% pero archivo independiente, sin dependencia cruzada entre cursos.
% ---------------------------------------------------------------------------

% ---- Paleta (alineada con las cajas del preámbulo: keybox/notebox/warnbox) ----
\definecolor{figblue}{HTML}{1F4E79}   % azul principal (L, curvas principales, keybox)
\definecolor{figred}{HTML}{B03A2E}    % rojo de énfasis (U, warnbox, error)
\definecolor{figamber}{HTML}{B8860B}  % ámbar (notebox, pivotes)
\definecolor{figgray}{HTML}{555555}   % auxiliares, ejes, ceros de matriz

% ---- TikZ: librerías y estilos reutilizables ----
% Nota: las puntas se declaran como `-{Latex[...]}` y no como `->`. La forma
% con llaves NO contiene un `>` literal, así que es inmune al `>` activo que
% introduce babel español (de todos modos se carga \usetikzlibrary{babel} en
% el bloque §6.3 para cubrir cualquier `->` suelto).
\usetikzlibrary{arrows.meta,calc,matrix,patterns}

\tikzset{
  figline/.style={figblue, line width=0.9pt},
  figaux/.style={figgray, dashed, line width=0.6pt},
  figaxis/.style={figgray, line width=0.7pt, -{Latex[length=2.2mm,width=1.5mm]}},
  % estructura de matrices (L triangular inferior, U triangular superior)
  matcelda/.style={minimum width=7mm, minimum height=7mm, anchor=center, font=\small},
  matnz/.style={matcelda, fill=figblue!12},
  matnzU/.style={matcelda, fill=figred!12},
  matcero/.style={matcelda, text=figgray!70},
  matpivote/.style={matcelda, fill=figamber!25, draw=figamber, line width=0.8pt},
  % nodos de interpolación / splines
  fignodo/.style={circle, inner sep=0pt, minimum size=4.5pt, draw=figblue, line width=0.9pt, fill=white},
  fignodoactivo/.style={fignodo, fill=figamber},
  % etiquetas
  figlbl/.style={font=\small},
  figlblaux/.style={font=\footnotesize, figgray},
}

% ---- pgfplots: estilos base ----
\pgfplotsset{
  % convergencia / error vs. n (grado, iteración) — suele ir en escala log
  errorfig/.style={
    width=10cm, height=6.5cm,
    xlabel style={font=\small},
    ylabel style={font=\small},
    tick label style={font=\footnotesize},
    every axis plot/.append style={line width=1.1pt, mark=*, mark size=1.5pt},
    grid=major, grid style={figgray!15},
    legend cell align=left,
    legend style={font=\footnotesize, draw=figgray!40, at={(0.5,-0.25)}, anchor=north},
    clip=false,
  },
  % interpolación / splines: un color por tramo, nodos marcados aparte
  interpfig/.style={
    width=10cm, height=6cm,
    axis lines=middle,
    xlabel style={font=\small, at={(axis description cs:1,0.5)}, anchor=west},
    ylabel style={font=\small, at={(axis description cs:0.5,1)}, anchor=south},
    tick label style={font=\footnotesize},
    every axis plot/.append style={line width=1.1pt},
    grid=none,
    legend cell align=left,
    legend style={font=\footnotesize, draw=figgray!40},
    clip=false,
  },
}

% (compilar desde el directorio de la clase para que la ruta relativa resuelva).
%
% Paleta propia de este curso (no se reusa la de Fisica3-2015/CDIV2017/
% ElecMag2024): mismos tonos base para alinear con keybox/notebox/warnbox,
% pero archivo independiente, sin dependencia cruzada entre cursos.
% ---------------------------------------------------------------------------

% ---- Paleta (alineada con las cajas del preámbulo: keybox/notebox/warnbox) ----
\definecolor{figblue}{HTML}{1F4E79}   % azul principal (L, curvas principales, keybox)
\definecolor{figred}{HTML}{B03A2E}    % rojo de énfasis (U, warnbox, error)
\definecolor{figamber}{HTML}{B8860B}  % ámbar (notebox, pivotes)
\definecolor{figgray}{HTML}{555555}   % auxiliares, ejes, ceros de matriz

% ---- TikZ: librerías y estilos reutilizables ----
% Nota: las puntas se declaran como `-{Latex[...]}` y no como `->`. La forma
% con llaves NO contiene un `>` literal, así que es inmune al `>` activo que
% introduce babel español (de todos modos se carga \usetikzlibrary{babel} en
% el bloque §6.3 para cubrir cualquier `->` suelto).
\usetikzlibrary{arrows.meta,calc,matrix,patterns}

\tikzset{
  figline/.style={figblue, line width=0.9pt},
  figaux/.style={figgray, dashed, line width=0.6pt},
  figaxis/.style={figgray, line width=0.7pt, -{Latex[length=2.2mm,width=1.5mm]}},
  % estructura de matrices (L triangular inferior, U triangular superior)
  matcelda/.style={minimum width=7mm, minimum height=7mm, anchor=center, font=\small},
  matnz/.style={matcelda, fill=figblue!12},
  matnzU/.style={matcelda, fill=figred!12},
  matcero/.style={matcelda, text=figgray!70},
  matpivote/.style={matcelda, fill=figamber!25, draw=figamber, line width=0.8pt},
  % nodos de interpolación / splines
  fignodo/.style={circle, inner sep=0pt, minimum size=4.5pt, draw=figblue, line width=0.9pt, fill=white},
  fignodoactivo/.style={fignodo, fill=figamber},
  % etiquetas
  figlbl/.style={font=\small},
  figlblaux/.style={font=\footnotesize, figgray},
}

% ---- pgfplots: estilos base ----
\pgfplotsset{
  % convergencia / error vs. n (grado, iteración) — suele ir en escala log
  errorfig/.style={
    width=10cm, height=6.5cm,
    xlabel style={font=\small},
    ylabel style={font=\small},
    tick label style={font=\footnotesize},
    every axis plot/.append style={line width=1.1pt, mark=*, mark size=1.5pt},
    grid=major, grid style={figgray!15},
    legend cell align=left,
    legend style={font=\footnotesize, draw=figgray!40, at={(0.5,-0.25)}, anchor=north},
    clip=false,
  },
  % interpolación / splines: un color por tramo, nodos marcados aparte
  interpfig/.style={
    width=10cm, height=6cm,
    axis lines=middle,
    xlabel style={font=\small, at={(axis description cs:1,0.5)}, anchor=west},
    ylabel style={font=\small, at={(axis description cs:0.5,1)}, anchor=south},
    tick label style={font=\footnotesize},
    every axis plot/.append style={line width=1.1pt},
    grid=none,
    legend cell align=left,
    legend style={font=\footnotesize, draw=figgray!40},
    clip=false,
  },
}

% (compilar desde el directorio de la clase para que la ruta relativa resuelva).
%
% Paleta propia de este curso (no se reusa la de Fisica3-2015/CDIV2017/
% ElecMag2024): mismos tonos base para alinear con keybox/notebox/warnbox,
% pero archivo independiente, sin dependencia cruzada entre cursos.
% ---------------------------------------------------------------------------

% ---- Paleta (alineada con las cajas del preámbulo: keybox/notebox/warnbox) ----
\definecolor{figblue}{HTML}{1F4E79}   % azul principal (L, curvas principales, keybox)
\definecolor{figred}{HTML}{B03A2E}    % rojo de énfasis (U, warnbox, error)
\definecolor{figamber}{HTML}{B8860B}  % ámbar (notebox, pivotes)
\definecolor{figgray}{HTML}{555555}   % auxiliares, ejes, ceros de matriz

% ---- TikZ: librerías y estilos reutilizables ----
% Nota: las puntas se declaran como `-{Latex[...]}` y no como `->`. La forma
% con llaves NO contiene un `>` literal, así que es inmune al `>` activo que
% introduce babel español (de todos modos se carga \usetikzlibrary{babel} en
% el bloque §6.3 para cubrir cualquier `->` suelto).
\usetikzlibrary{arrows.meta,calc,matrix,patterns}

\tikzset{
  figline/.style={figblue, line width=0.9pt},
  figaux/.style={figgray, dashed, line width=0.6pt},
  figaxis/.style={figgray, line width=0.7pt, -{Latex[length=2.2mm,width=1.5mm]}},
  % estructura de matrices (L triangular inferior, U triangular superior)
  matcelda/.style={minimum width=7mm, minimum height=7mm, anchor=center, font=\small},
  matnz/.style={matcelda, fill=figblue!12},
  matnzU/.style={matcelda, fill=figred!12},
  matcero/.style={matcelda, text=figgray!70},
  matpivote/.style={matcelda, fill=figamber!25, draw=figamber, line width=0.8pt},
  % nodos de interpolación / splines
  fignodo/.style={circle, inner sep=0pt, minimum size=4.5pt, draw=figblue, line width=0.9pt, fill=white},
  fignodoactivo/.style={fignodo, fill=figamber},
  % etiquetas
  figlbl/.style={font=\small},
  figlblaux/.style={font=\footnotesize, figgray},
}

% ---- pgfplots: estilos base ----
\pgfplotsset{
  % convergencia / error vs. n (grado, iteración) — suele ir en escala log
  errorfig/.style={
    width=10cm, height=6.5cm,
    xlabel style={font=\small},
    ylabel style={font=\small},
    tick label style={font=\footnotesize},
    every axis plot/.append style={line width=1.1pt, mark=*, mark size=1.5pt},
    grid=major, grid style={figgray!15},
    legend cell align=left,
    legend style={font=\footnotesize, draw=figgray!40, at={(0.5,-0.25)}, anchor=north},
    clip=false,
  },
  % interpolación / splines: un color por tramo, nodos marcados aparte
  interpfig/.style={
    width=10cm, height=6cm,
    axis lines=middle,
    xlabel style={font=\small, at={(axis description cs:1,0.5)}, anchor=west},
    ylabel style={font=\small, at={(axis description cs:0.5,1)}, anchor=south},
    tick label style={font=\footnotesize},
    every axis plot/.append style={line width=1.1pt},
    grid=none,
    legend cell align=left,
    legend style={font=\footnotesize, draw=figgray!40},
    clip=false,
  },
}
